\section{Provenance and Canonical Sources}

The hardened draft relies on the final claim inventory and the final paper bundle as its reporting boundary.

\paragraph{Canonical sources.}
All claim and figure sources used by the final draft are repo-local, completed, and canonical. In practice this means:
\begin{itemize}
\item final claim wording comes from \artifact{docs/research/11\_final\_claim\_inventory.md};
\item figure and table sources come from \artifact{docs/research/evidence/final\_paper\_20260403/};
\item historical interpretation comes from the backbone research docs and dated synthesis logs.
\end{itemize}

\paragraph{Explicit exclusions.}
The draft does not rely on evidence from:
\begin{itemize}
\item incomplete or superseded campaign roots,
\item scratch-only experiment directories,
\item raw Phase 3 roots moved under \artifact{/tmp/.../phase3\_raw},
\item optional aligned-refresh or confirmation reruns that were intentionally skipped by gate.
\end{itemize}

\paragraph{Why the final claim set is narrower than the project history.}
The project explored more configurations, revisions, and diagnostic paths than appear in the final main text. That narrowing is intentional. The final paper is not meant to re-narrate every operational step. It is meant to present the strongest bounded story supported by the canonical repo-local evidence. The dated logs remain the right place to preserve the full chronology and iteration history for handoff and future follow-up work.
